\section{A universal rank-detecting list}
\label{sec:rank}

The hitting list and its complexity bounds are now established.
We use the pencil-hitting consequence to recover the exact
noncommutative rank of a rectangular affine pencil from one list
that is independent of all coefficient values. The additional input
is the characterization of free-field rank by normalized
matrix-evaluation ranks, recalled in the proof.

\begin{corollary}
\label{cor:rank-detecting-list}
Let $a,b\geq1$ and $k\geq0$, and put $P=a+b+k+1$. From these
parameters in unary, a deterministic generator constructs a nonempty
finite list $\mathcal R_{a,b,k}$ of pairs
\[
 (e,X),\qquad e\geq1,\quad
 X=(X_1,\ldots,X_k)\in\Mat_e(\Q)^k,
\]
in a fixed polynomial in $P$ bit operations. Its number of entries,
total binary encoding length, and every matrix dimension $e$ are
bounded by fixed polynomials in $P$. The list depends only on $a,b,k$.

For every rectangular rational affine pencil
\[
 A(x)=A_0+\sum_{i=1}^k A_i x_i,\qquad A_i\in\Q^{a\times b},
\]
write $\operatorname{ncrk}A$ for its rank over the free skew field on
$x_1,\ldots,x_k$ over $\Q$, with ordinary rational rank when $k=0$.
Then the exact identity
\begin{equation}
 \operatorname{ncrk}A
 =\max_{(e,X)\in\mathcal R_{a,b,k}}
   \frac{\rank_{\Q}\!\left(A_0\otimes I_e+
             \sum_{i=1}^k A_i\otimes X_i\right)}{e}
 \label{eq:rank-detecting-list}
\end{equation}
holds for all choices and bit lengths of the coefficients.

Suppose that, at each requested pair $(e,X)$, an oracle returns
either the complete exact rational $ae$ by $be$ matrix in
\eqref{eq:rank-detecting-list}, in binary rational encoding, or its
exact ordinary rank. Querying the whole list nonadaptively and taking
the maximum of the exact quotients returns $\operatorname{ncrk}A$.
If $L$ is the total binary length of the responses, list generation,
query writing, and response processing use $\operatorname{poly}(P+L)$
bit operations, in addition to the sum of the oracle's costs.
The response lengths and oracle costs may depend on the coefficient
heights. If $a=0$ or $b=0$, the rank is zero and no queries are needed.
\end{corollary}

\begin{proof}
\emph{Affine evaluations and free-field rank.}
Put
$H(z)=\sum_{i=0}^k A_i z_i$, with a new variable $z_0$.
The coefficient-span invariance of Fortin and
Reutenauer~\citep[Section~2, Corollary~1, and Section~3, Lemma~1]{FR04}
includes the constant coefficient: the free-field rank of a
rectangular affine linear matrix depends only on the span of all its
coefficient matrices. The coefficient spans of $A$ and $H$ are both
$\Span_{\Q}\{A_0,\ldots,A_k\}$, so their free-field ranks agree.
This uses coefficient-span invariance, not a specialization $z_0=1$
in a free skew field.

For $e\geq1$, let $h_e$ be the maximum rank of
$\sum_{i=0}^k A_i\otimes Z_i$ over $Z_i\in\Mat_e(\Q)$, and let
$\alpha_e$ be the maximum rank of
$A_0\otimes I_e+\sum_{i=1}^k A_i\otimes X_i$ over rational $X_i$.
Zero-padding the $A_i$ to square matrices of order $\max(a,b)$
preserves both free-field rank and every ordinary blow-up rank.
The homogeneous blow-up characterization and its attainment over
$\Q$~\citep[arXiv version, Definition~1.2 and its discussion,
Theorem~1.5]{IQS18} therefore give
\[
 r:=\operatorname{ncrk}A=\max_{e\geq1}\frac{h_e}{e},
\]
with the maximum attained. The tensor order used in that reference
differs from $A_i\otimes Z_i$ only by row and column permutations.
Here padding is used to preserve ranks, not to test invertibility.

For each fixed $e$, taking $Z_0=I_e$ gives $\alpha_e\leq h_e$.
If $h_e>0$, some $h_e$ by $h_e$ minor is a nonzero polynomial $p$
in the entries of the $Z_i$. The polynomial $p\det Z_0$ is nonzero;
since $\Q$ is infinite, it is nonzero at a rational tuple. There
$Z_0$ is invertible and the rank is $h_e$. Right multiplication gives
\[
 \left(\sum_{i=0}^k A_i\otimes Z_i\right)
       (I_b\otimes Z_0^{-1})
 =A_0\otimes I_e+\sum_{i=1}^k A_i\otimes(Z_iZ_0^{-1}),
\]
without changing rank. The case $h_e=0$ is immediate, so
$\alpha_e=h_e$ for every $e$. Consequently
\begin{equation}
 r=\max_{e\geq1}\frac{\alpha_e}{e}
 \label{eq:affine-normalized-rank}
\end{equation}
is attained: every individual affine evaluation has rank at most
$re$, and some rational evaluation has rank exactly $re$. No bound
on the dimension or entry heights of this existence witness
will be used.

\emph{Square completions for rank thresholds.}
For $0\leq t\leq\min(a,b)$, form the square affine pencil
\[
 C_t=\begin{pmatrix}A&U_t\\ V_t&0\end{pmatrix},
 \qquad U_t:a\times(a-t),\quad V_t:(b-t)\times b.
\]
The entries of $U_t,V_t$ are independent variables, so their
dimension-$e$ evaluations may be arbitrary matrices of the displayed
block sizes multiplied by $e$. For a fixed rational tuple $X$,
there exist rational choices of these completion blocks making
$C_t(X,U_t,V_t)$ invertible if and only if
$\rank A(X)\geq te$. Necessity follows by splitting the
completed matrix into its $A(X)$ and two off-diagonal blocks:
\[
 (a+b-t)e=\rank C_t(X,U_t,V_t)
 \leq\rank A(X)+(a-t)e+(b-t)e.
\]
For sufficiency, choose a $te$-dimensional subspace
$K\subseteq\Q^{be}$ on which $A(X)$ is injective. Take $V_t$ to be
surjective with kernel $K$, and take $U_t$ to map isomorphically onto
a complement of $A(X)K$ in $\Q^{ae}$. A vector in the kernel of the
completed matrix then has its first component in $K$, and the direct
sum $A(X)K\oplus\operatorname{im}U_t$ makes both components zero.
This proves invertibility. It includes $t=0$ and all empty-block
cases; in particular $C_0$ can always be made invertible.

\emph{A coefficient-independent list.}
Reserve $a^2$ variables $u_{ij}$ and $b^2$ variables $v_{ij}$,
disjoint from the first $k$ variables. For $C_t$, use the first
$a-t$ columns of the $u$-pool for $U_t$ and the first $b-t$ rows of
the $v$-pool for $V_t$. The pools can be reused for different $t$.
Thus every $C_t$ is a pencil in the same
$N=k+a^2+b^2$ variables, with unused variables allowed, and has
order $a+b-t\leq a+b$. Set $s=a+b$, so this order is at most
$B=2s+1$. Both $N$ and $s$ are positive. Apply the single list
$\mathcal H_{N,s}$ from Corollary~\ref{cor:pencil-hitting}, and
project each output onto its first $k$ matrices, retaining its
dimension $e$ and any repetitions. This defines $\mathcal R_{a,b,k}$;
it is nonempty because $C_0$ is feasible.

For $t=r$, equation~\eqref{eq:affine-normalized-rank} and the
completion criterion make $C_r$ feasible. A hit for $C_r$ therefore
has projected rank at least $re$, while every projected evaluation
has rank at most $re$. This proves~\eqref{eq:rank-detecting-list}
exactly, without rounding; individual ranks need not be divisible by
$e$. The argument includes $r=0$. When $k=0$, the projected tuple is
empty but still carries $e$, and the identity also follows directly
from $\rank(A_0\otimes I_e)=e\rank A_0$.

Since $N+s+1$ is polynomial in $P$, Proposition~\ref{prop:generator-complexity}
and the dimension bound~\eqref{eq:generator-dimension-bounds} give
the stated generation, query-encoding, and dimension bounds. The
generator never constructs $C_t$, learns $r$, or reads the
coefficients: $C_r$ is used only in the existence proof. For complete
matrix responses, exact rational elimination is polynomial in the
encoded response lengths and dimensions. For ordinary-rank responses,
only exact comparison of the rank/dimension quotients is needed.
These observations give the stated cost accounting.
\end{proof}

The assertion concerns exact matrix values or ordinary ranks, not
singularity-bit access to the original pencil. It makes no claim for
arbitrary rational circuits, finite fields, or total oracle-computation
bit time independent of coefficient heights. If complete scalar value
queries at zero and the standard basis points are also available, then
$k+1$ such queries reveal all coefficients, after which known
coefficient-aware rank algorithms apply with those bit lengths
counted~\citep[arXiv version, Theorem~1.5]{IQS18}. The distinct
conclusion here is the universal nonadaptive coefficient-independent
list, which also works with ordinary-rank responses.

